
	
\section{Algorithme de Bröker, Charles et Lauter}

\subsection{Enoncé et contexte}

On souhaite utiliser la théorie précédente pour établir une procédure de calcul d'isogénie plus rapide que la méthode naïve donnée par la formule de Vélu. Plus précisément on suppose connaître 
\begin{enumerate}
	\item Un modèle de courbe elliptique $E$ ordinaire, donnée par une équation de Weierstrass réduite définie sur un corps $\F_{q}$ de caractéristique $p$. On précalcule le nombre de point du modèle $|E(\F_{q})|$.
	\item L'anneau d'endomorphisme $End(E)\simeq\OR_{\Delta}$ donné par son discriminant $\Delta$. On connaît donc le corps quadratique imaginaire $\K$ et le conducteur $f$. On supposera de plus que l'isomorphisme $End(E)\simeq\OR_{\Delta}$ est normalisé. 
	\item Un nombre premier $l\neq p$.
	\item Un entier $n\in\N$ et un point $P\in E(\F_{q^{n}})$. 
\end{enumerate} 

Et on cherche une courbe elliptique $E'$ définie sur $\F_{q}$ telle qu'il existe une $l$-isogénie séparable $\phi : E\rightarrow E'$ définie sur $\F_{q}$. De plus on souhaite pouvoir calculer $\phi(P)$. 

\begin{center}
	\textbf{Dans toutes la suite, on suppose que $l$ ne divise pas $pf_{m}|E(\F_{q^{n}})|$, que $l$ est décomposé dans $End(E)$ et que $f_{m}$ et $|E(\F_{q^{n}})|$ sont premier entre eux}
\end{center}

Cette condition nous oblige à éviter certains nombres premiers. Elle a cependant plusieurs avantages :
\begin{itemize}
	\item Sous cette condition, $l$ ne divise pas $f_{m}$ donc toutes les $l$-isogénies séparables sont horizontales. Comme $l$ est décomposé, il existe $2$ $l$-isogénies et leurs noyaux sont accessibles comme groupes de torsions d'idéaux de norme $l$. 
	\item On peut même représenter une isogénie uniquement par l'idéal qui engendre son noyau. Un tel idéal peut être donné sous la forme $\left[ l, \frac{c + d\pi_{q}}{f_{m}}\right]$ donc seulement déterminé par $l$, $c$ et $d$, trois entier qui en pratique seront de $\log(l)$ bits. On peut représenter une $l$-isogénie en $O(log(l))$ bits. 
	\item Si on suppose seulement que $l$ ne divise pas $f$, il existe des isogénies descendantes que l'on ne peut pas approcher avec la théorie précédente. Tandis que si $l$ ne divise pas $f_{m}$, on peut s’intéresser à toutes les $l$-isogénies définies sur $\F_{q}$.
	\item Si $l$ ne divise pas $pf_{m}|E(\F_{q^{n}})|$, $f_{m}$ et $|E(\F_{q^{n}})|$ sont inversibles modulo $l$. Le fait que $f_{m}$ soit inversible sera utile par la suite pour manipuler des endomorphisme de la forme $\left[ \frac{c + d\pi_{q}}{f_{m}}\right] $. Le fait que $|E(\F_{q^{n}})|$ soit inversible sera utile pour normaliser l'isogénie obtenue (Voir section 5.5, 5.6).
	\item La dernière condition est cruciale pour la dernière étape (Voir section 5.6 proposition 5.29). On suppose avoir vérifié lors du choix de la courbe elliptique $E$ que cette condition est vérifié pour les valeurs de $n$ qui nous intéresse.
	
\end{itemize}

On déduit une procédure naïve à partir de la formule de Vélu, à condition de savoir expliciter le générateur d'un groupe de torsion.

\begin{algorithm}
	\caption{Calcul d'une forme explicite normalisée de $l$-isogénie}
	\begin{algorithmic}[1]
		\Require Un modèle défini sur $\F_{q}$ de courbe elliptique $E$ à multiplication complexe, le discriminant $\Delta$ de $End(E)$, un nombre $l$ premier ne divisant pas $p*f_{m}*|E(\F_{q^{n}})|$ et tel que $\left( \frac{\Delta}{l}\right) = 1$
		\Ensure Une forme explicite normalisée d'une $l$-isogénie $\phi : E \rightarrow E'$ partant de $E$.
		\State Factoriser $l\OR_{\Delta}$ avec le théorème de Kummer
		\State Choisir un des deux idéaux $\LR$ de norme $l$ obtenu
		\State Calculer un générateur $P$ du groupe de torsion $E\left[ \LR\right]$ (voir section 5.4)
		\State Appliquer la formule de Vélu pour trouver une forme explicite $\phi_1 : E \rightarrow E'$ 
		\State Calculer l'invariant différentiel de $E'$
		\State Calculer la constante $c$ telle que $\phi_{1}^{*}(w_{E_{1}}) = cw_E$. 
		\State Appliquer l'isomorphisme $\mu_{c} : (x,y) \mapsto (c^{-2}x, c^{-3}y)$ à la forme explicite de $\phi_1 : E\rightarrow E'$ pour calculer $\phi = \mu_{c}\circ\phi_{1} : E \rightarrow E'$
		\State \Return La forme explicite $\phi : E \rightarrow E'$
	\end{algorithmic}
	
\end{algorithm}


Cependant la formule de Vélu demande de parcourir le noyau entier : la procédure est linéaire en $l$, ce qui n'est pas satisfaisant. L'algorithme de Broker, Charle et Lauter consiste à améliorer l'algorithme 1 en se ramenant à des idéaux de petites normes grâce à l'action du groupe de classe $Cl(\OR_{\Delta})$ sur $Ell_{\F_{q}}(\OR_{\Delta})$. On présente les différnetes étapes de la façon suivante :

\begin{itemize}
	\item On a besoin d'outils pour calculer l'action de $Cl(\OR_{\Delta})$: on étudie cela dans la section 5.2. On utilise des formes quadratiques pour manipuler les classes d'idéaux. On explique comment trouver une famille génératrice de $Cl(\OR_{\Delta})$ représentée par des idéaux de petites normes premier bien choisies. On explique ensuite comment exprimer une classe donnée dans cette famille génératrice.
	\item On aura tout de même besoin de la formule de Vélu au cours de l'algorithme. Dans la section 5.3, on étudie comment l'appliquer efficacement, et on analyse la complexité obtenue. 
	\item Pour appliquer la formule de Vélu, il nous faut savoir trouver un générateur d'un groupe de torsion d'un idéal donné. Dans la section 5.4 on explique comment en obtenir un lorsque l'idéal est de norme premier.
	\item Partant d'un idéal $\LR$ de norme $l$, on pourra exprimer sa classe $\CL$ dans une famille génératrice bien choisie. On en déduit une expression de $\LR$ à un idéal fractionnaire principal prêt qu'il faut expliciter. On explique comment trouver un générateur d'un idéal principal dans la section 5.5.
	\item La procédure de Broker Charle Lauter sera énoncé en section 5.6, puis la complexité sera étudié en section 5.7. 
\end{itemize}







\subsection{Calculs dans les groupes de classes}


Pour mener des calculs dans un groupe de classe $Cl\left(\OR_{\Delta} \right)$, il nous faut d'abord résoudre deux problèmes :
\begin{enumerate}
	\item Pour manipuler une classe sans ambiguïté il faut définir un représentant normal, unique dans chaque classe. En particulier on souhaite pouvoir tester rapidement si deux idéaux sont dans la même classe.
	\item On voudra calculer des isogénie associée aux idéaux représentant les générateurs d'une famille génératrice. Il faut dont que leurs normes vérifient les conditions fixées sur $l$ : On veut pouvoir choisir des représentants de classes de sorte que leurs normes soient des nombres premiers bien choisis. 
\end{enumerate}

Un outil va nous permettre de résoudre les deux problèmes en même temps : les formes quadratiques à coefficient entier. 



\textbf{Déplacer la partie théorique des formes quadratiques}


\subsubsection{Forme standard des idéaux et formes quadratiques binaires}

On va étudier une correspondance entre d'une part les idéaux d'un ordre d'un corps quadratique imaginaire $\OR_{\Delta}$, et d'autre part des formes quadratiques définies positives de discriminant $\Delta$. On commence par définir les formes quadratiques qui nous intéressent :

\begin{Def}
	
	Soit $f = ax^2 + bxy + cy^2$ une forme quadratique. $f$ est binaire si elle est définie positive et à coefficient dans $\Z$. Son discriminant est $\Delta = b^2 - 4ac$. On notera $f = (a,b,c)$ pour désigner une forme quadratique binaire, ou $f = (a,b)$ si le discriminant $\Delta$ est connu. $f$ est dite primitive si ses coefficients $a$, $b$ et $c$ sont premiers dans leur ensemble.
\end{Def}

On peut associer à une forme quadratique binaire primitive de discriminant $\Delta$ un idéal propre de $\OR_{\Delta}$.

\begin{Prop}\cite{Coh00}[Th 7.7]
	
	Soit $f = (a,b)$ une forme quadratique binaire primitive de discriminant $\Delta$. L'application :
	\begin{equation*}
		I : (a,b)\mapsto\left[ a, \frac{-b + \sqrt{\Delta}}{2}\right] 
	\end{equation*}
	associe à $f$ un idéal propre de $\OR_{\Delta}$.
\end{Prop}

\begin{Def}
	Un idéal entier $\LR\subset\OR_{\Delta}$ est dit primitif s'il est de la forme $I(a,b)$. 
\end{Def}

Pour définir la réciproque de l'application $I$, il faut pouvoir retrouver les coefficients $a$ et $b$ à partir d'une écriture adéquate de $\LR$ :

\begin{Prop}
	
	Tout idéal entier $\LR\subset\OR_{\Delta}$ s'écrit 
	\begin{equation*}
		\LR = m\left[ a , \frac{-b + \sqrt{\Delta}}{2}\right] 
	\end{equation*}
	avec $m,a \in\N$, $b\in\Z$, $b^{2} = \Delta [4a]$, et $c = (b^2 -\Delta)/4a$ avec $a, b, c$ premiers entre eux. L'écriture est unique avec $b$ défini modulo $2a$ : on peut choisir $-a<b\leq a$.
\end{Prop}

\begin{proof}\cite{Buchmann95}[2.5]
	La preuve citée est rédigée dans le cas des corps quadratique réel, c'est à dire de discriminant positif. Mais à part un choix de convention différent pour fixer $b$, la preuve s'adapte sans difficulté au cas imaginaire.
\end{proof}

\begin{Def}
	Un idéal $\LR$ est sous forme standard s'il est donné par l'unique base définie par la proposition 5.2. Lorsque $\Delta$ est connue, on notera $\LR = (m,a,b)$ une écriture standard de $\LR$. On dira que $\LR$ est primitif si $\LR = (1,a,b)$.
	
\end{Def}

\begin{Rem}
	On peut en déduire une écriture standard d'un idéal fractionnaire $\IR$ de $\OR_{\Delta}$. En effet d'après la proposition 1.11 il existe $d\in\N$ minimal tel que $d\IR = \LR$ soit entier, de forme standard $\LR = (m,a,b)$. Alors la forme standard de $\IR$ est
	\begin{equation*}
		\LR = \frac{m}{d}\left[ a , \frac{-b + \sqrt{\Delta}}{2}\right]
	\end{equation*}
\end{Rem}


\begin{Coro}
	Soit $\LR = (m,a,b)$ un idéal entier de $\OR_{\Delta}$ sous forme standard. Soit $c = \frac{\Delta - (b)^2}{4a}$. L'application
	\begin{equation*}
		F : (m,a,b)\mapsto (a,b) = ax^2 + bxy + cy^2
	\end{equation*}
	associe à $\LR$ une forme quadratique binaire de même discriminant. 
	
\end{Coro}

\begin{Prop}
	Les applications $I$ et $F$ induisent une bijection entre les idéaux entiers primitifs et les formes quadratiques binaires primitives.
\end{Prop}

\begin{proof}
	Par construction, si $\LR = (1,a,b)$ est un idéal entier primitif de $\OR_{\Delta}$, alors $I(F(\LR)) = \LR$ et $F(I((a,b))) = (a,b)$.
\end{proof}

\begin{Rem}
	On peut de même associé une forme quadratique à un idéal fractionnaire à partir de sa forme standard. 
\end{Rem}

On peut retrouver facilement la norme et l'inverse d'un idéal entier donné sous forme standard.

\begin{Prop}\cite{Coh00}[5.2.5]
	
	Soit $\LR = (m,a,b)\subset\OR_{\Delta}$ un idéal entier sous forme standard. Soit $\LR' = \left( \frac{m}{N(\LR)},a,-b\right) $ (qui est donc un idéal fractionnaire de $\OR_{\Delta}$). Alors
	\begin{enumerate}
		\item $N(\LR) = m^2a$
		\item $\LR\LR' = \left[ m, \frac{-b + \sqrt{\Delta}}{2}\right] $
		\item $\LR$ est propre si et seulement s'il est primitif, et dans ce cas $\LR' = \LR^{-1}$
	\end{enumerate}
\end{Prop}

On cherche à trouver un représentant d'un classe d'idéaux de $\OR_{\Delta}$ de façon systématique et pratique pour l'implémentation. Prendre un idéal primitif ne suffit pas : Il peut y avoir plusieurs idéaux primitifs dans la même classe. 

\begin{Def}
	Deux formes quadratiques binaires primitives sont équivalentes si leurs idéaux primitifs associés sont dans la même classe. On note $Cl(\Delta)$ les classes d'équivalences de formes quadratiques binaires primitives de discriminant $\Delta$.
\end{Def}

La bijection précédente entre les idéaux entiers primitif et les formes quadratiques binaires primitive induit une structure de groupe sur les formes quadratiques binaires primitives. L'isomorphisme passe au quotient et induit un isomorphisme entre $Cl(\OR_{\Delta})$ et $Cl(\Delta)$. Cependant la notion d'équivalence entre formes quadratiques binaires primitives peut être définie sans faire le lien avec les idéaux de $\OR_{\Delta}$, ainsi qu'une loi de groupe sur les formes et sur les classes via le produit le Dirichlet \cite{Cox}[3.8]

\begin{Def}
	Soit $f = (a,b)$, $f' = (a', b')$ deux formes quadratiques binaires primitives de même discriminant $\Delta$. On défini le produit de Dirichlet $f*f'$ de la façon suivante : Soit $B$ défini modulo $2aa'$ par
	\begin{enumerate}
		\item $B = b [2a]$
		\item $B = b' [2a']$
		\item $B^2 = \Delta [4aa']$
	\end{enumerate}
	Alors $f*f' = (aa', B)$ est une forme quadratique binaire primitive.
\end{Def}

\begin{Prop}\cite{Cox}[Th 7.7]
	
	Avec les notation précédentes, $I(f*g) = I(f)I(g)$ et $F(\LR\LR') = F(\LR)*F(\LR')$ Et ces opérations induisent des isomorphismes entre groupes de classes.
\end{Prop}




\textbf{Fin partie théorique à déplacer}




On défini une notion de forme quadratique binaire réduite, qui aura l'avantage d'être unique dans chaque classe :

\textbf{Expliquer rapidement l'action des matrices, utile plus tard aussi}

\begin{Def}
	Soit $f = (a,b)$ une forme quadratique binaire primitive sous forme standard de discriminant $\Delta < 0$. Soit $c = (\Delta - b^2)/4a$. $f$ est une forme quadratique réduite si $|b| \leq a \leq c$ et si de plus lorsque $a = c$, $b \geq 0$. Un idéal entier primitif $\LR$ de $\OR_{\Delta}$ est dit réduit si $F(\LR)$ est une forme quadratique réduite
\end{Def}

\begin{Prop}\cite{Cox}[2.8]
	
	Soit $\OR_{\Delta}$ un ordre d'un corps quadratique imaginaire. Toute classe de $Cl(\Delta)$ contient une unique forme quadratique réduite. Toute classe de $Cl(\OR_{\Delta})$ contient un unique idéal entier primitif réduit. 
\end{Prop}

\begin{Rem}
	Une forme quadratique réduite $f = (a,b)$ vérifie 
	\begin{equation*}
		|b|\leq|a|\leq\sqrt{\frac{|\Delta|}{3}}
	\end{equation*}
	On peut donc représenter chaque classe $Cl(\OR_{\Delta})$ en $O(\log(|\Delta|))$ bits. 
\end{Rem}

La preuve citée de la proposition 5.7 fournit un algorithme donnant l'élément réduit de la classe en $\OR(\log(\Delta)^{2} )$ à partir d'un représentant quelconque.

\begin{algorithm}
	\caption{Réduction de forme quadratique binaire primitive}
	\begin{algorithmic}[1]
		\Require Une forme quadratique binaire primitive $(a,b,c)$
		\Ensure Une forme quadratique réduite équivalente.
		\State $test \gets False$
		\While{test = False}
		\If{ $|b| > a$}
		\State Calculer $q$, $r$ tels que $b = 2aq + r$, $-a<r\leq a$
		\State $c \gets c - bq + aq^2$
		\State $b \gets r$
		\EndIf
		\State $test \gets True$
		\If{$a > c$}
		\State Echanger $a$ et $c$
		\State $b \gets -b$
		\State $test \gets False$
		\EndIf
		\If{$b < 0$}
		\If{$a = -b$ OU $a = c$}
		\State $b \gets -b$
		\EndIf
		\EndIf
		\EndWhile
		\State \Return $(a,b,c)$
	\end{algorithmic}
\end{algorithm}



\begin{Rem}
	On peut donc vérifier que deux idéaux entiers sont dans la même classe : il suffit de réduire leurs formes quadratiques binaires primitives associées et vérifier que l'on trouve le même résultat. 
\end{Rem}






\subsubsection{Représentant de classe premier avec le conducteur}

La forme réduite sera adapté pour faire des calculs dans le groupe de classe. Pour le calcul d'isogénie, on voudra cependant trouver un système de représentant de normes bien choisies. A priori l'idéal associé à la forme réduite ne convient pas. 

\begin{Def}
	Un entier $m\in\N$ est représenté par une forme quadratique binaire primitive $f$ s'il existe $s$, $t$ deux entiers tels que $m = f(s,t)$.
\end{Def}

\begin{Prop}\cite{Cox}[Th 7.7]
	
	Un entier $m\in\N$ est représenté par une forme quadratique binaire primitive $f$ si et seulement s'il existe un idéal $\LR\in \widetilde{I(f)}$ tel que $N(\LR) = m$
\end{Prop}

\begin{Rem}
	La preuve de cette proposition fournit une méthode pour trouver l'idéal $\LR$ de norme $m$ à partir de deux entiers $s$, $t$ tel que $f(s,t) = 1$ \cite{Cox}[Lem 2.3]. Cependant si $m$ est premier, on connaît les idéaux de normes $m$ lorsque $m$ est décomposé ou ramifié: Un idéal de norme $m$ contient $m\OR_{\Delta}$ et est donc donné par le théorème de Kummer. Ceci permet de trouver l'idéal $\LR$ sans expliciter les entiers $s$ et $t$.
\end{Rem}

\begin{Prop}\cite{Cox}[Cor 7.17]
	
	Soit $M\in\N$ en entier. Toutes formes quadratiques binaires primitives représente au moins un entier premier avec $M$. Ainsi toutes classes d'idéaux d'un ordre quadratique imaginaire possède un représentant premier avec $M$. 
	
\end{Prop}

Ainsi, toute classe d'idéaux possède un représentant premier avec $M = pf_{m}|E(\F_{q^{n}})|$. Comme il y a deux choix d'idéaux de norme $p$ premier lorsque $p$ est décomposé, on va fixer une convention nous permettant de choisir. 

\begin{Def}
	Soit $p$ premier décomposé dans $\OR_{\Delta}$. On fixe un algorithme prenant entrée $p$ et $\OR_{\Delta}$ qui calcule une racine de $\OR_{\Delta}$ modulo $4p$. La forme quadratique première de discriminant $\Delta$ associée à $p$ est $(p,b_{p})$ où $b_{p}\in\N$ vérifie $b_{p}^2 = \Delta [4p]$ et est obtenue par l’algorithme choisit. L'idéal premier associé à $p$ est alors $\PR_{p} = I((p,b_{p}))$. En particulier $N(\PR_{p}) = p$.
\end{Def}

\begin{Rem}
	L'idéal premier associé à $p$ est donc l'un des deux idéaux premier qui factorise $p$ dans $\OR_{\Delta}$. $I((p,-b_{p}))$ est le second, et représente la classe de $(\PR_{p})^{-1}$ d'après la proposition 5.5. 
\end{Rem}












\subsubsection{Famille génératrice du groupe de classe}




On veut former une famille génératrice de $Cl(\OR_{\Delta})$ définie avec un système de représentant formé uniquement d'idéaux de norme $p$ adaptée au calcul d'isogénie, c'est à dire $p$ premier avec $pf_{m}|E(\F_{q^{n}})|$. Une idée naïve est d'établir la liste croissante des nombres premiers acceptables plus petit qu'une borne $N$. Pour $N$ suffisamment grand on peut espérer avoir assez de représentant de classes distinctes parmi les idéaux premier associés pour obtenir une famille génératrice.

\begin{Def}
	Soit $N\in\N$. On note $P_{N} = \left\lbrace \PR_{p} : p \; premier, 0\leq p\leq N, \left( \frac{\Delta}{p}\right) = 1\right\rbrace$ la liste des idéaux premiers associés aux nombres premiers plus petits que $N$. 
\end{Def}

La proposition suivante permet de choisir la borne $N$.

\begin{Prop}\cite{Seysen}[Th 3.3]
	
	Sous l'hypothèse de Riemann généralisée, il existe une constante $C$ telle que $Cl(\OR_{\Delta})$ soit engendré par les classes des éléments de $P_{N}$, avec 
	\begin{equation*}
		N = C\times \log(\Delta)^{2}
	\end{equation*}
	
\end{Prop}

\begin{Rem}
	Pour estimer la constante $C$, \cite{Seysen} cite deux résultats proposant $C= 280$ ou même $C = 2$. \cite{Coh00} propose $C = 6$. Cependant d'après \cite{Coh00}, en pratique cette borne $N$ est pessimiste et l'on conjecture une borne en $O(\log(\Delta)^{1+\epsilon})$ pour tout $\epsilon > 0$.
\end{Rem}

\begin{Rem}
	Comme les nombres premiers divisant $pf_{m}|E(\F_{q^{n}})|$ sont en nombre fini, les exclure de $P_{N}$ ne change pas le comportement asymptotique de la borne $N$. Prenant en compte ceci et la remarque précédente, on suppose donc que $P_{N}$ reste une famille génératrice si l'on ignore les diviseurs de $pf_{m}|E(\F_{q^{n}})|$. 
\end{Rem} 

Comme on peut choisir $N$ de taille asymptotiquement raisonnable, on peut appliquer la procédure naïve, détaillée dans l'algorithme 3. 


\begin{algorithm}
	\caption{Famille génératrice du groupe de classe}
	\begin{algorithmic}[1]
		\Require Un entier positif $N$, un discriminant $\Delta$, la caractéristique $p$, le conducteur maximal $f_{m}$, le cardinal $|E(\F_{q^{n}})|$.
		\Ensure Une liste $L_{premiers}$ de nombres premiers plus petit que $N$ ne divisant pas $pf_{m}|E(\F_{q^{n}})|$, la liste $L_{ideaux}$ de leurs idéaux de Seysen respectif, la liste $L_{formes}$ de leurs formes de Seysen respectives. 
		\State $L_{premiers}$, $L_{ideaux}$, $L_{formes} \gets$ \O
		\State $n \gets 2$
		\While{$n < N$}
		\If{$\left( \frac{\Delta}{n}\right) = 1$ ET $n$ ne divise pas $pf_{m}|E(\F_{q^{n}})|$}
		\State Mettre $n$ dans $L_{premiers}$
		\State Calculer $q_{n} = (n ,b_{n})$ la forme quadratique première associée à $n$ de discriminant $\Delta$.
		\State Mettre $q_{n}$ dans $L_{formes}$
		\State Calculer $I_{n}$ l'idéal premier associé à $n$. 
		\State Mettre $I_{n}$ dans $L_{ideaux}$
		\EndIf
		\State $n \gets nextprime(n)$
		\EndWhile
		\State\Return $L_{premiers}$, $L_{ideaux}$, $L_{formes}$
	\end{algorithmic}
\end{algorithm}









\subsubsection{Écriture d'une classe dans une famille génératrice}


Soit $\LR\subset\OR_{\Delta}$ un idéal entier. Soit $L_{premiers}$, $L_{ideaux}$ et $L_{formes}$ les listes retournées par l'algorithme $3$. On note $k$ leurs tailles. Alors il existe $e_{1},\ldots,e_{k}\in\Z$ tels que
\begin{equation*}
	\CL = \CL_{1}^{e_{1}}\CL_{2}^{e_{2}}\ldots\CL_{k}^{e_{k}}
\end{equation*}
avec $\LR_{1}\ldots\LR_{k}\in L_{premiers}$. De plus il existe $\alpha\in\K$ tel que
\begin{equation*}
	\LR = (\alpha\OR_{\Delta})\LR_{1}^{e_{1}}\LR_{2}^{e_{2}}\ldots\LR_{g}^{e_{g}}
\end{equation*}
Dans cette section on ne cherche pas à expliciter $\alpha$ mais seulement des exposants $e_{i}$. Une idée naïve consiste à chercher de tels exposants au hasard. Deux faits vont permettre de rendre cela effectif :

\begin{enumerate}
	\item Certaines classes représentées par des formes réduites ayant de bonnes propriétés vont être explicitement et rapidement factorisable.
	\item Les marches aléatoires dans le groupe de classe vont avoir de bonnes probabilités de terminer en un temps sous-exponentiel sur une classe facilement factorisable.
\end{enumerate}

D'abord on cite le théorème sur la factorisation de classes:

\begin{The}\cite{Seysen}[Th 3.1]
	
	Soit $\LR\subset\OR_{\Delta}$ un idéal tel que $\CL$ ait une forme réduite $(a,b)$. Si $a = p_{1}^{e_{1}}p_{2}^{e_{2}}\ldots p_{g}^{e_{g}}$ est la décomposition en facteurs premiers de l'entier $a$, alors
	\begin{enumerate}
		\item $p_{i}$ est décomposé dans $\OR_{\Delta}$.
		\item Si $\PR_{p_{i}} = (1,p_{i},b_{p_{i}})$ est l'idéal de Seysen de $p_{i}$, alors $b = \pm b_{p_{i}} $ mod $2a$.
		\item $\CL = \CP_{p_{1}}^{\pm e_{1}}\CP_{p_{2}}^{\pm e_{2}}\ldots\CP_{p_{g}}^{\pm e_{g}}$ 
		\item L'exposant de $\PR_{p_{i}}$ est $+e_{i}$ si et seulement si $b = b_{i} [2a]$
	\end{enumerate}
\end{The}

\begin{Coro}
	Avec les notations du théorème précédent, si $a$ se factorise entièrement avec des premiers $p_{i}\in L_{premiers}$ obtenue par l'algorithme 3, alors $\CL$ est facilement factorisable connaissant la factorisation de $a$. 
\end{Coro}

On énonce ensuite les bonnes propriétés des marches aléatoires sur les groupes de classes. Sans rappeler la théorie des graphes expanseurs, on cite les résultats principaux.

\begin{Def}
	Soit $S$ une partie génératrice d'un groupe $G$. Le graphe de Cayley de $G$ selon $S$ est défini ainsi :
	\begin{enumerate}
		\item Les sommets sont associés chacun à un élément de G
		\item Une arêtes relie deux sommets associés à $g$ et $g'$ si et seulement s'il existe $s\in S$ tel que $g = sg'$
	\end{enumerate}
	S est supposé stable par inverse pour que le graphe ne soit pas orienté.
\end{Def}

\begin{Prop}\cite{Jao_2009}[Lem 2.1]
	
	Soit $G = Cl(\OR_{\Delta})$ et $S = L_{ideaux}\cup (L_{ideaux})^{-1}$. On considère une marche aléatoire en $t$ étape partant d'un sommet quelconque du graphe de Cayley de G selon S. Si 
	\begin{equation*}
		t = O(\frac{\log(|\Delta|)}{\log(\log(|\Delta|))})
	\end{equation*} 
	Alors pour toute partie $S_{0}\subset G$, la marche aléatoire termine dans $S_{0}$ avec probabilité au moins $\frac{|S_{0}|}{2|G|}$. 
\end{Prop}

Concrètement, dès que l'on enchaîne des marches aléatoires de taille de l'ordre de $\frac{\log(|\Delta|)}{\log(\log(|\Delta|))}$, les résultats seront uniformément répartis parmi les éléments de $Cl(\OR_{\Delta})$.

\begin{Rem}
	Puisque le groupe est abélien, la norme $l1$ du vecteurs des exposants, que l'on notera $x$, est au pire le nombre de pas $t$ de la marche aléatoire. 
\end{Rem}

\begin{Coro}
	On peut choisir des vecteurs d'exposants aléatoire parmi les vecteurs $x$ tels que 
	\begin{equation*}
		\left\| x\right\| _{1} \leq C\times\frac{\log(|\Delta|)}{\log(\log(|\Delta|))}
	\end{equation*}
	
\end{Coro}

\textbf{Estimation de la constante C ?}

\begin{algorithm}
	\caption{Factorisation de la classe d'un idéal dans une famille génératrice donnée}
	\begin{algorithmic}[1]
		\Require Un entier positif $N$, une borne de pas de marche aléatoire $t$, un discriminant $\Delta$, la caractéristique $p$, le conducteur maximal $f_{m}$, le cardinal $|E(\F_{q^{n}})|$. Un idéal $\LR\subset\OR_{\Delta}$. 
		\Ensure Deux listes $L_{exposants}$, $L_{facteurs}$ décrivant une factorisation de $\CL$.
		\State $L_{ideaux}$, $L_{formes}$, $L_{premiers}\gets$ résultat de l'algorithme 3 appliqué à $N$, $\Delta$, $p$, $f_{m}$, $|E(\F_{q^{n}})|$. 
		\State $classe_{test}\gets$ forme réduite de $\LR$
		\State $k\gets$ taille de $L_{ideaux}$
		\State $L_{x}\gets$ liste de taille $k$ initialisée à $0$
		
		\While{$classe_{test}$ ne se factorise pas avec le théorème 5.12}
		\State $L_{x}\gets$ liste de taille $k$ initialisée à $0$
		\For{$0<i<t $}
		\State $Indice\gets$ indice aléatoire de $L_{x}$.
		\State $Increment\gets\pm1$
		\State $L_{x}\left[ Indice\right] \gets L_{x}\left[ Indice\right] + Increment$
		\EndFor
		\For{$0\leq i \leq k - 1$}
		\State $x_{i}\gets L_{x}[i]$
		\For{$1 \leq j \leq x_{i}$}
		\State $classe_{test}\gets classe_{test}L_{formes}[i]$
		\State $classe_{test}\gets classe_{test}$ réduit par l'algorithme 2
		\EndFor
		\EndFor
		\EndWhile
		
		\State $L_{y}\gets$ une liste de taille $k$ contenant les exposants la factorisation de $classe_{test}$ donnée par le théorème 5.12.
		\State $L_{exposants}\gets L_{y} - L_{x}$
		\State $L_{facteurs}\gets L_{ideaux}$
		\For{$e_i$ parcourant $L_{exposants}$ }
		\If{$e_i < 0$}
		\State $e_i\gets-e_i$
		\State $\left( L_{facteurs}\right) _{i} \gets $ idéal conjugué de $\left( L_{facteurs}\right) _{i}$
		\EndIf
		\EndFor
		
		\State \Return $L_{exposants}$, $L_{facteurs}$
	\end{algorithmic}
\end{algorithm}

 





\subsection{Algorithme SquareVélu}

Étant donné un modèle de courbe elliptique $E$ défini sur $\F_{q}$, les formules de Vélu permettent de construire une forme explicite d'une isogénie $\phi$ à partir de la connaissance de son noyau. Pour une $l$-isogénie, le noyau peut être donné sous deux formes : un point $K\in E(\F_{q^{r}})$ générateur du noyau, avec $r \geq 1$, ou un polynôme $F$ dont les racines sont les abscisses des points finis du noyau. Dans le second cas, il faut adapter les formules pour les écrire en fonction de $F$.

\begin{Def}
	Soit $E$ un modèle de courbe elliptique et $\phi : E\rightarrow E'$ une $l$-isogénie, avec $l$ premier impair. On sépare $Ker(\phi)$ en une partition $\left\lbrace \OR \right\rbrace \cup R\cup R'$ où $R$ regroupe $(l-1)/2$ points finis distincts, et $R'$ leurs inverses. Un polynôme $F$ décrit le noyau de $\phi$ sur $E$ si c'est un polynôme unitaire de degré $(l-1)/2$ dont les racines sont les abscisses des points finis de $R$. Un tel polynôme est unique. 
\end{Def}

Dans les deux cas, \cite{Velusqrt} propose une méthode pour évaluer des polynômes de degré $l$ dans $\F_{q}[X]$ en $O(\sqrt{l})$ opérations dans $\F_{q}$. Les résultats suivants en découle.

\begin{Prop}
	Soit $E$ un modèle de courbe elliptique défini sur $\F_{q}$, et $\phi : E\rightarrow E'$ une $l$-isogénie, avec $l$ premier impair. Soit $K$ un générateur de $Ker(\phi)$, et $P\in E(\F_{q})$. Il est possible de calculer un modèle de $E'$ et $\phi(P)\in E'$ en $O(\sqrt{l})$ opérations dans $\F_{q}$, en utilisant l'algorithme de \cite{Velusqrt} que l'on appelera SquareVélu.
\end{Prop}

\begin{Rem}
	SquareVélu est une amélioration asymptotique de l'application naïve des formules de Vélus. Cependant pour les petits nombres premiers, l'application naïve reste plus rapide. En pratique avec l’implémentation de Sage, il est conseillé d'utiliser SquareVélu uniquement si $l > 100$. 
\end{Rem}

On décrit maintenant comment adapter les formules de Vélu à l'utilisation de polynômes décrivant le noyau d'une isogénie.

\begin{Prop}\cite{Kohel}[\S 2.4]
	
	Soit $E : y^2 = x^3 + ax + b$ un modèle de courbe elliptique défini sur $\F_{q}$, et $\phi$ une $l$-isogénie, avec $l$ premier impair. Soit $F$ le polynôme décrivant le noyau de $\phi$ sur $E$. Alors on peut obtenir une forme explicite de $\phi$ avec
	\begin{equation*}
		\phi(x,y) = \left( \frac{G(X)}{F(X)^2}\;,\;\frac{H(X,Y)}{F(X)^3}\right) 
	\end{equation*}
	où $G$ et $H$ peuvent être obtenue explicitement à partir de $F$. Pour $P\in E(\F_{q})$, on peut alors calculer $\phi(P)$ en $O(\sqrt{l})$ opérations dans $\F_{q}$ via les méthodes de \cite{Velusqrt}
\end{Prop}

\begin{Rem}
	Il existe aussi des formules adaptées au cas où $l = 2$
\end{Rem}










\subsection{Générateur d'un groupe de torsion}





\subsubsection{Méthode générale}

Dans cette partie, on suppose connu un idéal $\LR\subset\OR_{\Delta}$ de norme $l$ premier décomposé, et une courbe elliptique $E$ définie sur $\F_{q}$ à multiplication complexe avec $End(E)\simeq\OR_{\Delta}$ normalisé. On cherche un générateur du groupe de torsion $E\left[ \LR\right]\subset E(\Fb_{p})$. On suppose $l$ premier avec le conducteur maximal $f_{m}$.

\begin{Rem}
	Si $\LR$ donné sous forme standard $\left[ l, \frac{b + \sqrt{d_{\K}}}{2}\right]$, $\LR$ est primitif car de norme premier. Il est facile de passer à une écriture dépendant de $w_{\K}$ au lieu de $\sqrt{d_{K}}$, et la proposition 4.23 permet alors d'obtenir en quelques calculs une base de la forme 
	\begin{equation*}
		\LR = \left[ l, \frac{c + d\pi_{q}}{f_{m}}\right]
	\end{equation*}
	On suppose $\LR$ donné sous cette forme. 
\end{Rem}

On sait alors que $E\left[ \LR\right]$ est par définition l'intersection de $E\left[ l\right] $ et du noyau de $\left[  \frac{c + d\pi_{q}}{f_{m}}\right]$. Dans $E[l]$, à partir d'une base $P$, $Q$, on peut obtenir les générateurs des $l+1$ noyaux de $l$-isogénie possible. (voir proposition 2.6). On veut alors pouvoir tester lequel est dans $E\left[  \frac{c + d\pi_{q}}{f_{m}}\right]$. 

\begin{Prop}
	Avec les notations précédentes, si $l$ ne divise pas $f_{m}$, alors les isogénies $ \left[  \frac{c + d\pi_{q}}{f_{m}}\right]$ et $\left[ c + d\pi_{q}\right] $ restreintes à $E\left[ l\right] $ ont même noyau.
	
\end{Prop}

\begin{proof}
	Puisque $f_{m}$ et $l$ sont premier entre eux, $\left[ f_{m} \right]$ restreinte à $E\left[ l \right]$ est bijective. Or l'application $\left[.\right] $ est un morphisme d'anneau, donc $\left[  \frac{c + d\pi_{q}}{f_{m}}\right] \circ \left[ f_{m }\right] =  \left[ c + d\pi_{q}\right] $, ce qui permet de conclure.
\end{proof}

On en déduit l’algorithme suivant, proposé dans \cite{BrokerCharlesLauter}[\S 3.2] comme méthode générale :

\begin{algorithm}
	\caption{Point générateur d'un groupe de $\LR$-torsion, méthode 1}
	\begin{algorithmic}[1]
		\Require Une courbe elliptique $E$ à multiplication complexe définie sur $\F_{q}$, d'anneau d'endomorphisme $\OR_{\Delta}$, un nombre $l$ premier ne divisant pas $pf_{m}|E(\F_{q^{n}})|$ et tel que $\left( \frac{\Delta}{l}\right) = 1$, un idéal propre $\LR\subset\OR_{\Delta}$ sous la forme $\LR = \left[ l, \frac{c + d\pi_{q}}{f_{m}}\right]$
		\Ensure Un point $K\in E(\Kb)$ générateur de $E[\LR]$. 
		\State Calculer le polynôme de $l$-division sur $E$.
		\State En déduire une extension $\F_{q^r}$ sur laquelle le polynôme de $l$-division se factorise. (Ainsi $E(\F_{q^r})$ contient $E[l]$).
		\State Calculer une base $P$, $Q$ de $E\left[ l\right] $. 
		\For{$i$ allant de $0$ à $l - 1$}
		\If{$\left[ c + d\pi_{q}\right](P + iQ) = 0$} 
		\Return la valeur $P + iQ$ 
		\EndIf
		\EndFor
		\State \Return $Q$
	\end{algorithmic}
\end{algorithm}

Cette algorithme peut être implémenté sur sage, mais pose plusieurs problème :

\begin{enumerate}
	\item L'extension $\F_{q^r}$ peut être très grande, Broker Charle et Lauter affirme que $r$ est en général de l'ordre de $l$. 
	\item Le polynôme de $l$-division est degré $\frac{l^2 - 1}{2}$ et peut être lourd à calculer, et à factoriser.
\end{enumerate}

En pratique sur l'exemple "small" de \cite{BrokerCharlesLauter}, $l = 31$ mais le calcul prend déjà une dizaine de secondes.







\subsubsection{Optimisation de la recherche d'extension de corps}

On propose une amélioration de la première méthode : le but est de trouver directement la plus petite extension qui contient la $\LR$-torsion.\\

On remarque d'abord que, partant de la forme standard de l'idéal $\LR$, on peut toujours se ramener à $d = 1$ (il suffit d'expliciter les calculs)\textbf{Si l'on retire la section 4 il faut détailler un peu ici}. Si $P$ est un générateur du groupe de torsion, alors $\pi_{q}(P) = -cP$ : $P$ est un vecteur propre de $\pi_{q}$ restreint à $E[l]$. On se ramène donc à une étude $\pi_{q}$ restreint à $E[l]$.

\begin{Prop}\cite{MMSTV}[Prop 1]
	
	Soit $E$ définie sur $\F_{q}$ de $j$-invariant différent de $0$, $1728$.
	\begin{enumerate}
		\item $E$ admet une unique $l$-isogénies définie sur $\F_{q}$ si et seulement si $\pi_{q}$ restreint à $E[l]$ a un unique espace propre de dimension $1$.
		\item $E$ admet exactement deux $l$-isogénies définies sur $\F_{q}$ si et seulement si $\pi_{q}$ restreint à $E[l]$ a deux valeurs propres distinctes .
		\item $E$ admet $l+1$ isogénies définies sur sur $\F_{q}$ si et seulement si $\pi_{q}$ restreint à $E[l]$ a une valeur propre double.
	\end{enumerate}
	
\end{Prop}

\textbf{ Cas de 0 isogénie ? } Puisque $l$ est supposé décomposé et premier avec $f_{m}$, on sait qu'il existe exactement $2$ isogénies définies sur $\F_{q}$. On note $v_{1}$ et $v_{2}$ les deux valeurs propres de $\pi_{q}$ restreint à $E[l]$. 

\begin{Prop}\cite{MMSTV}
	
	Soit $G_{1}$, $G_{2}$ les deux sous groupes de $E$ noyaux des deux $l$-isogénies définies sur $\F_{q}$. Soit $r_{1}$, $r_{2}$ minimaux tels que $G_{i}\subset E\left( \F_{q^{r_{i}}}\right) $. Alors $r_{i}$ est l'ordre de $v_{i}$ dans $\F_{l}^{*}$, et $r_{i}$ divise $l-1$. De plus si $r_{i}$ est pair, les abscisses des points de $G_{i}$ sont dans $\F_{q^{r_{i}/2}}$.
\end{Prop}

On sait donc directement dans quelle extension chercher les générateurs des noyaux. On note $r$ son degré. On cherche alors $\K\in E(\F_{q^{r}})$ un vecteurs propre de $\pi_{q}$ pour la valeurs propres $-c$, d'ordre $l$. L'idée et que si la structure de groupe de $E(F_{q^{r}})$ est 
\begin{equation*}
	E(\F_{q^{r}}) \simeq \Z/n_{1}\Z\times\Z/n_{2}\Z
\end{equation*}

avec $n_{1}$ divisant $n_{2}$, alors un point choisit au hasard sera probablement d'ordre $n_{2}$. Cependant on connait seulement le cardinal $N = n_{1}n_{2}$

\begin{Prop}\cite{Sil09}[Exercice 5.6]
	
	Avec les notations précédentes, $n_{1}^2$ divise $N$ et $q^r = 1$ modulo $n_{1}$.
\end{Prop}

On en déduit l'algorithme 6 :
\begin{algorithm}
	\caption{Point générateur d'un groupe de $\LR$-torsion, méthode 2}
	\begin{algorithmic}[1]
		\Require Une courbe elliptique $E$ à multiplication complexe définie sur $\F_{q}$, d'anneau d'endomorphisme $\OR_{\Delta}$, un nombre $l$ premier ne divisant pas $pf_{m}|E(\F_{q^{n}})|$ et tel que $\left( \frac{\Delta}{l}\right) = 1$, un idéal propre $\LR\subset\OR_{\Delta}$ sous la forme $\LR = \left[ l, \frac{c + d\pi_{q}}{f_{m}}\right]$
		\Ensure Un point $K\in E(\Kb)$ générateur de $E[\LR]$. 
		\State $r\gets$ ordre de $-c$ dans $\F_{l}^{*}$.
		\State $N\gets$ Cardinal de $E(\F_{q^{r}})$
		\State $n_{max} \gets pgcd(N, q^r - 1)$
		\State $L\gets$\O
		\For{$i$ divisant $n_{max}$}
		\If{$i^2$ divise $N$}
		\State Mettre $i$ dans $L$
		\EndIf
		\EndFor
		\State $P\gets$ Point aléatoire de $E(\F_{q^{r}})$
		\For{$i$ parcourant $L$}
		\State $K\gets (N/i)P$
		\If{$K\neq \OR$}
		\If{$\pi_{q}(K)=-cK$}
		\Return $K$ 
		\EndIf
		\EndIf
		\EndFor
		\State Reprendre l'étape $10$
	\end{algorithmic}
\end{algorithm}

	
\subsubsection{Méthode d'Elkies}

La méthode 2 précédente est une net amélioration de la méthode 1, mais nécessite parfois de manipuler des groupes $E(F_{q^{r}})$ de très grand cardinal, lorsque $r$ est de l'ordre de $l-1$. En pratique, cela arrive suffisamment souvent pour être problématique. On présente ici une troisième méthode, proposé par \cite{Bostan_2008}[\S 4.3, fastElkies'] ayant pour avantage d'éviter tout calcul dans des extensions de $\F_{q}$.\\ 

On résume l'idée de la méthode en admettant les résultats utilisés. Soit $E$ une courbe elliptique à multiplication complexe de $j$-invariant $j$, définie sur $\F_{q}$.
\begin{itemize}
	
	\item Il existe un polynôme $\Psi_{l}(X,Y)$, dépendant uniquement de $l$, tel que les racines de $\Psi_{l}(j,Y)$ dans $\Fb_{p}$ sont exactement les $j$-invariants des $l+1$ courbes $l$-isogènes à $E$. Si les racines sont dans $\F_{q}$, elles correspondent aux courbes images des $l$-isogénies définies sur $\F_{q}$
	
	\item Soit $h$ une racine de $\Psi_{l}(j,Y)$ dans $\F_{q}$, alors il est possible d'obtenir un modèle défini sur $\F_{q}$ de la courbe $E'$ de $j$-invariant $h$, à partir de la donnée de $E$, $h$ et $l$.
	
	\item  Soit $p_{1}\in\Qb$ la somme des abscisses des $l-1$ points finis du noyau de la $l$-isogénie $\phi : E\rightarrow E'$. Alors $p_{1}\in\F_{q}$ peut être obtenues à partir de la connaissance de $E$, $h$ et $l$, en menant des calculs uniquement dans $\F_{q}$.
	
	\item Soit $F_{h}$ le polynôme dont les racines sont les abscisses des élément du noyau de $\phi$. Soit $F'_{h}$ défini de la même façon à partir d'un relèvement de During de $E$, $E'$ $\phi$. Soit $\wp$, $\wp'$ les fonction de Weierstrass associée au relèvement de $E$ et $E'$. Alors
	\begin{equation}
		z^{l-1}F'_{h}(\wp(z)) = \exp(S)
	\end{equation}
	où $S$ est une série entière dont tout les coefficients dépendent de l'écriture en série de Laurent de $\wp$, $\wp'$ et du relèvement de $p_{1}$. Les coefficients sont des éléments de $\Qb$
	
	\item On peut alors retrouver les coefficients de $F'_{h}$ en identifiants les $(l-1)/2$ premiers termes des deux membres de l'équation. Si chaque terme a une bonne réduction modulo l'idéal de réduction $\PR$ au dessus de $p$, alors on obtient directement les coefficient de $F_{h}$. La bonne réduction n'est cependant pas garanti si l'on doit calculer trop de termes : il faut que $l$ soit 'petit', plus précisément
	\begin{equation*}
		l < 4*(p+1)
	\end{equation*}
\end{itemize}

Cette méthode a été proposé par Elkies, et amélioré dans \cite{Bostan_2008}. Les améliorations visent à minimiser et optimiser les calculs de séries.

\begin{Def}
	En suivant la convention utilisée par Sage, on appellera BMSS l'algorithme proposé dans \cite{Bostan_2008}[\S 4.3, fastElkies'], prenant en entrée $E$, $E'$ deux modèles de courbe elliptique à multiplication complexe définie sur $\F_{q}$ et $l < 4*(p+1)$ un nombre premier tel qu'il existe au moins une $l$-isogénie horizontale normalisée $\phi : E \rightarrow E'$. Cette algorithme renvoie un polynôme $F$ tels que les abscisses des points finis du noyaux de $\phi$ sont les racines du polynôme $F$ (sans redondance).  
\end{Def}

\begin{Prop}
	Soit $M : \N\times\N\rightarrow\N$ une fonction telle que $M(n,q)$ soit la complexité de la multiplication de deux polynômes de degrés au plus $n$ sur $\F_{q}$. L'algorithme BMSS admet une complexité en $O\left( M(l)log(l)\right)$
\end{Prop}

Lorsque $l$ est décomposé, $\Psi_{l}(j,Y)$ admet exactement deux racines dans $\F_{q}$. Si l'on souhaite trouvé l'unique isogénie normalisée de noyau $E[\LR]$ où $\LR$ est un idéal de norme $l$, on ne peut pas connaître à l'avance le bon choix de $j$-invariant. il faut donc un moyen de tester si le premier choix était bon, et sinon recommencer avec la seconde racine. On veut toujours éviter de calculer dans une extension de $\F_{q}$

\begin{Prop}
	Soit $E : y^2 = x^3 + ax + b$ un modèle de courbe elliptique à multiplication complexe définie sur $\F_{q}$, d'anneau d'endomorphisme $\OR_{\Delta}$. Soit $\LR = \left[ l, \frac{c + \pi_{q}}{f_{m}}\right]$ un idéal de $\OR_{\Delta}$ de norme $l$ premier. Soit $F$ un polynôme sur $\F_{q}$.  On pose 
	\begin{equation*}
		R = \F_{q}\left[ X,Y\right] / \left( F, Y^2 - X^3 - aX - b\right) 
	\end{equation*}
	Soit $(R_{1}(X), YR_{2}(X))$ les fonctions rationnelles définissant l'endomorphisme $[-c]$. Si les racines de $F$ dans $\Fb_{p}$ sont les abscisses des points finis de $E[\LR]$, alors dans le corps de fraction de $R$,
	\begin{equation*}
		\bar{X}^q = R_{1}(\bar{X})\;,\;\bar{Y}^{q} = \bar{Y}R_{2}(\bar{X})
	\end{equation*}
\end{Prop}

On en déduit l'algorithme $7$.

\begin{algorithm}
	\caption{Méthode d'Elkies}
	\begin{algorithmic}[1]
		\Require Une courbe elliptique $E$ à multiplication complexe définie sur $\F_{q}$, d'anneau d'endomorphisme $\OR_{\Delta}$, un nombre $l$ premier ne divisant pas $pf_{m}|E(\F_{q^{n}})|$ et tel que $\left( \frac{\Delta}{l}\right) = 1$, un idéal propre $\LR\subset\OR_{\Delta}$ sous la forme $\LR = \left[ l, \frac{c + d\pi_{q}}{f_{m}}\right]$, le polynôme modulaire $\Psi_{l}$.
		\Ensure Un polynôme $F\in\F_{q}\left[ X\right]$ dont les racines sont les abscisses des points finis de $E[\LR]$, sans redondance. 
		\State $P_{invariants}(Y)\gets PGCD(\Psi_{l}(j,Y),Y^q - Y)$ ($P_{invariants}$ est de degré 2)
		\State $h_{1}, h_{2}\gets$ racines de $P_{invariants}$
		\State $E_{1}\gets$ modèle normalisé de $h_{1}$ (formule explicite)
		\State $F_{1}\gets BMSS(E,E_{1},l)$
		\If{$F_{1}$ vérifie la proposition 5.24}
		\State \Return $F_{1}$
		\Else 
		\State $E_{2}\gets$ modèle normalisé de $h_{2}$ (formule explicite)
		\State $F_{2}\gets BMSS(E,E_{2},l)$
		\State \Return $F_{2}$
		\EndIf
	\end{algorithmic}
\end{algorithm}






\subsection{Problème de l'Idéal Principal}

Dans cette partie, on suppose connaître une factorisation de la classe $\CL$ obtenue avec l'algorithme 4
\begin{equation}
	\CL = \CP_{p_{1}}^{e_{1}}\CP_{p_{2}}^{e_{2}}\ldots\CP_{p_{g}}^{e_{g}}
\end{equation}
Où les idéaux $\PR_{p_{i}}$ sont des idéaux de Seysen de normes $p_{i}$ différent de $p$ et ne divisant pas $f_{m}|E(\F_{q^{n}})|$. On cherche alors $\alpha\in\K$ tel que 
\begin{equation}
	\LR = (\alpha\OR_{\Delta})\PR_{p_{1}}^{e_{1}}\PR_{p_{2}}^{e_{2}}\ldots\PR_{p_{g}}^{e_{g}}
\end{equation}
On commence par se ramener à chercher un élément de $\OR_{\Delta}$ : Comme $\PR_{i} = I_{p_{i}} = (p_{i},b_{i})$, la forme $(p_{i},-b_{i})$ représente $\CP_{i}^{-1}$ et $\bar{\PR_{i}} = I((p_{i},-b_{i}))$. Donc
\begin{equation}
	\LR\bar{\PR_{1}}^{e_{1}}\bar{\PR_{2}}^{e_{2}}\ldots\bar{\PR_{g}}^{e_{g}} = \beta\OR_{\Delta}
\end{equation}
où $\beta\OR_{\Delta}$ est un idéal entier. Ainsi $\beta\in\OR_{\Delta}$.

\begin{Prop}
	Avec les notations précédentes, Soit $m = p_{1}^{e_{1}}p_{2}^{e_{2}}\ldots p_{g}^{e_{g}}$. Alors $\alpha = \frac{\beta}{m}$.
\end{Prop}

\begin{proof}
	Partant de l'équation $(2)$, on multiplie chaque membre par $\LR\bar{\PR_{1}}^{e_{1}}\bar{\PR_{2}}^{e_{2}}\ldots\bar{\PR_{g}}^{e_{g}}$. Le membre de gauche est égale à $\beta\OR_{\Delta}$ par définition de $\beta$ par l'équation $(3)$. Le membre de gauche est
	\begin{equation*}
		(\alpha\OR_{\Delta})\PR_{p_{1}}^{e_{1}}\PR_{p_{2}}^{e_{2}}\ldots\PR_{p_{g}}^{e_{g}}\bar{\PR_{1}}^{e_{1}}\bar{\PR_{2}}^{e_{2}}\ldots\bar{\PR_{g}}^{e_{g}}
	\end{equation*}
	Or d'après la proposition 1.17.3, $\PR_{p_{i}}^{e_{i}}\bar{\PR_{i}}^{e_{i}} = p_{i}^{e_{i}}\OR_{\Delta}$, donc 
	\begin{equation*}
		(\alpha\OR_{\Delta})\PR_{p_{1}}^{e_{1}}\PR_{p_{2}}^{e_{2}}\ldots\PR_{p_{g}}^{e_{g}}\bar{\PR_{1}}^{e_{1}}\bar{\PR_{2}}^{e_{2}}\ldots\bar{\PR_{g}}^{e_{g}} = \alpha m\OR_{\Delta}
	\end{equation*}
	Donc $\beta = \alpha m$
\end{proof}

On cherche plutôt à obtenir $\beta$. Puisque l'on connaît l'idéal $\LR$ et les idéaux $\PR_{p_{i}}$ sous forme standard, on connaît la forme standard de $\beta\OR_{\Delta}$ grâce à l'équation $(3)$. De plus le produit d'idéaux primitif est primitif (C'est une conséquence de l'isomorphisme entre les idéaux primitif et les formes quadratiques binaires primitives). On se ramène donc au problème suivant : étant donné un idéal entier primitif principal $(1,a,b)$, on doit trouver un générateur $\beta$. Comme $N((1,a,b)) = a = N_{\Q}^{\K}(\beta)$, connaître $a$ donne la norme de $\beta$. On sait que $\beta = \frac{x + y \sqrt{\Delta}}{2}$ avec $x$, $y$ entiers.

\begin{Prop}
	
	Soit $\LR = [a , \frac{-b + \sqrt{\Delta}}{2}]$ un idéal de $\OR_{\Delta}$ primitif sous forme standard. Soit $f = F(\LR)$ sa forme quadratique binaire primitive associée. Alors 
	\begin{equation*}
		f(x,y) = ax^2 + bxy + cy^2 = N_{\Q}^{\K}\left( ax - y\frac{-b + \sqrt{\Delta}}{2}\right) / N(\LR)
	\end{equation*}
\end{Prop}

\begin{proof}
	On sait que par définition de $N_{\Q}^{\K}$, 
	
	\begin{align*}
		N_{\Q}^{\K}\left( ax - y\frac{-b + \sqrt{\Delta}}{2}\right) / N(\LR) &= \left( ax + y\frac{b}{2} + y\frac{\sqrt{\Delta}}{2}\right) \left( ax + y\frac{b}{2} - y\frac{\sqrt{\Delta}}{2}\right)/ a\\
		&= \left( \left( ax + y\frac{b}{2} \right)^{2} - \left( y\frac{\sqrt{\Delta}}{2}\right)^{2}\right)/a\\
		&=  ax^2 + bxy + y^2\left( \frac{b^2 - \Delta }{4a}\right) \\
		&= ax^2 + bxy + cy^2 = f(x,y)
	\end{align*}
	
	la dernière ligne étant dû au fait que $f$ est de discriminant $\Delta = b^2 - 4ac$.
	
\end{proof}


\begin{Coro}
	Soit $\LR = (1,a,b)\subset\OR_{\Delta}$ un idéal primitif sous forme standard. $\LR$ est principal si et seulement si $1$ est représenté par $f$, c'est à dire s'il existe deux entiers $x, y\in\Z$ tels que $f(x,y) = 1$. De plus dans ce cas,
	\begin{equation*}
		\LR = \left( ax - \frac{-b + \sqrt{\Delta}}{2}\right) \OR_{\Delta}
	\end{equation*}
\end{Coro}

\begin{proof}
	$\LR$ est un idéal principal si et seulement s'il existe un élément de $\LR$ de norme égale à $N(\LR)$. Puisque $\LR = [a , \frac{-b + \sqrt{\Delta}}{2}]$, tout élément $\beta$ de $\LR$ est de la forme
	\begin{equation*}
		\beta = \left( ax - \frac{-b + \sqrt{\Delta}}{2}\right)
	\end{equation*}
	Alors $N_{\Q}^{\K}(\beta) = N(\LR) \iff f(x,y) = 1$ d'après ce qui précède.
\end{proof}

Dans le contexte de notre problème, on sait que l'idéal est principal par construction, et on connaît la forme quadratique $f$ associée. On cherche $(x,y)\in\Z^2$ solution de $f(x,y) = 1$. La notion de forme réduite va nous permettre de conclure 

\begin{Prop}\cite{Cox}[Thm 2.8, 2.9]
	
	Soit $r = ax^2 + bxy + cy^2$ une forme quadratique binaire primitive. Si $r$ est réduite, alors tout les entiers représentés par $r$ sont positifs et $a$ est le plus petit entier représenté par $r$, avec $r(1,0) = a$.
\end{Prop}

Ainsi lorsque $f$ est une forme quadratique binaire primitive qui représente $1$, sa forme réduite $r$ est de la forme $x^2 \pm xy + cy^2$. 

\begin{Rem}
	Chaque étape de l'algorithme 2 correspond à un changement de variable en $(x,y)$ a appliquer à la forme quadratique \cite{Cox}[Thm 2.8]. En effet
	\begin{enumerate}
		\item $(x,y)\mapsto (x - 2qy, y)$ correspond aux lignes $5$ et $6$.
		\item $(x,y)\mapsto (x + y, y)$ correspond aux lignes $17$ et $18$ dans le cas $a = -b$.
		\item $(x,y)\mapsto(-y,-x)$ correspond aux lignes $10$ et $11$, ainsi qu'aux lignes $17$ et $18$ dans le cas $a = c$. 
	\end{enumerate}
\end{Rem}

On en déduit un algorithme pour trouver le générateur $\beta$ d'un idéal principal primitif.

\begin{algorithm}
	\caption{Générateur d'un idéal principal primitif donné sous forme standard}
	\begin{algorithmic}[1]
		\Require $(1,a,b)$ un idéal primitif d'un ordre $\OR_{\Delta}$ dont on sait qu'il est principal.
		\Ensure Un élément $\beta$ tel que $(1,a,b) = \beta\OR_{\Delta}$. 
		\State $f\gets ax^2 + bxy + cy^2$
		\State $r\gets$ résultat de l'algorithme 2 appliqué à $f$
		\State $L\gets$ liste des changements de variables utilisés durant la réduction de $f$
		\State $(x,y)\gets(1,0)$
		\State $(x,y)\gets(x',y')$ obtenus en appliquant tout les changements de variables retenus dans $L$ dans le sens inverse.
		\State \Return $(x,y)$
	\end{algorithmic}
\end{algorithm}









\subsection{Application au calcul d'isogénie}

On se donne un modèle $E$ de courbe elliptique à multiplication complexe définie sur $\F_{q}$, et $\LR$ un idéal de norme $l$ vérifiant toutes les conditions voulues.

\begin{center}
	\textbf{On rappelle que $l$ est supposé décomposé dans $End(E)$, ne divise pas $pf_{m}|E(\F_{q^{n}})|$ et que $f_{m}$ et $|E(\F_{q^{n}})|$ sont premier entre eux}
\end{center}

On est capables de factoriser la classe $\CL$ en un produit de classes représentées par des idéaux $\PR_{1}, \ldots,\PR_{N}$ de petites normes vérifiant toutes les mêmes conditions que $l$. On peut donc calculer des isogénies dont les noyaux sont des groupes de $\PR_{i}$-torsion, et ainsi reproduire l'action de $\CL$ sur $E$ en suivant sa factorisation. 

\begin{Def}
	Soit $\phi : E \rightarrow E/E[\LR]$. Soit $\CL = \CP_{1}^{e_{1}} \ldots\CP_{N}^{e_{N}}$ une factorisation de $\CL$. Les composantes de $\phi$ associée à cette factorisation sont isogénies obtenues par la suite
	\begin{equation*}
		E\rightarrow E/E\left[ \PR_{1}\right] \rightarrow E/E\left[ \PR_{1}^{2}\right] \rightarrow\ldots\rightarrow E/E\left[ \PR_{1}^{e_{1}}\right] \rightarrow
	\end{equation*}
	\begin{equation*}
		\rightarrow E/E\left[ \PR_{1}^{e_{1}}\PR_{2}\right] \rightarrow\ldots\rightarrow E/E\left[ \PR_{1}^{e_{1}}\PR_{2}^{e_{2}}\right] \rightarrow\ldots\rightarrow E/E\left[\PR_{1}^{e_{1}} \ldots\PR_{N}^{e_{N}}\right]
	\end{equation*}
	Les composantes explicites sont les uniques formes explicites normalisées des composantes de sorte que la première parte du modèle $E$. On note $E_{c}$ le modèle explicite de $E/E\left[\PR_{1}^{e_{1}} \ldots\PR_{N}^{e_{N}}\right]$ obtenue avec le calcul de la dernière composante. On note $\phi_c : E\rightarrow E_c$ la composée de toutes les composantes
\end{Def}

$E_c$ est un modèle de $E/E[\LR]$. Cependant la $l$-isogénie $\phi : E \rightarrow E/E[\LR]$ n'est pas égale à $\phi_c$, qui n'est pas de degré $l$ sauf exception. On a cependant la relation
\begin{equation*}
	\LR = (\alpha\OR_{\Delta})\PR_{p_{1}}^{e_{1}}\PR_{p_{2}}^{e_{2}}\ldots\PR_{p_{N}}^{e_{N}}
\end{equation*}
où $\alpha\OR_{\Delta}$ est un idéal principal fractionnaire. On sait calculer $\alpha = \beta/m$ avec $\beta = \frac{u + v\pi_{q}}{f_{m}}\in\OR_{\Delta}$. On a choisi les $p_{i}$ premier avec $|E(\F_{q^{n}})|$, donc $m$ est aussi premier avec ce cardinal. \textbf{On a de plus supposé que $f_{m}$ est premier avec $|E(\F_{q^{n}})|$}, ce qui permet de calculer $k = (mf_{m})^{-1}$ modulo $|E(\F_{q^{n}})|$. L'endomorphisme $k(u + v\pi_{q})$ est donc bien défini. 

\begin{Prop}
	Avec les notations précédentes, 
	\begin{equation*}
		\forall\; P\in E[|E(\F_{q^{n}})|]\;,\;\phi(P) = (k(u + v\pi_{q}))\circ\phi_c(P)
	\end{equation*}
\end{Prop}

\begin{proof}
	Soit $\tilde{E}$, $\tilde{E_c}$, $\tilde{E'} = \widetilde{E/E\left[\LR\right] }$ trois courbe elliptique définie sur $\Qb$, obtenue par le théorème de relèvement de Deuring appliquée via l'endomorphisme de Froebenius, et via le même idéal de réduction $\PR$ au dessus de la caractéristique $p$. Soit $\Lambda$ un réseau de $\C$ tel que $E_{Lambda} = \tilde{E}$. Alors la factorisation de $\LR$ implique que
	\begin{equation*}
		\tilde{E'}\simeq\C/\LR^{-1}\Lambda = \C/\left( (\alpha\OR_{\Delta})\PR_{p_{1}}^{e_{1}}\PR_{p_{2}}^{e_{2}}\ldots\PR_{p_{N}}^{e_{N}} \right)^{-1}\Lambda \simeq \tilde{E'}
	\end{equation*}
	De plus,
	\begin{equation*}
		\tilde{E_c} \simeq \C/\left(\PR_{p_{1}}^{e_{1}}\PR_{p_{2}}^{e_{2}}\ldots\PR_{p_{N}}^{e_{N}} \right)^{-1}\Lambda
	\end{equation*}
	Par définition, $\phi_{\LR}$ et $\phi_c$ correspondent au niveau des tores au morphisme identité (voir section 3.4). L'action de $\alpha$ correspond donc à l'identité. On pose
	\begin{equation*}
		\phi_{\alpha} : \C/\left(\PR_{p_{1}}^{e_{1}}\PR_{p_{2}}^{e_{2}}\ldots\PR_{p_{N}}^{e_{N}} \right)^{-1}\Lambda \rightarrow \C/\left( (\alpha\OR_{\Delta})\PR_{p_{1}}^{e_{1}}\PR_{p_{2}}^{e_{2}}\ldots\PR_{p_{N}}^{e_{N}} \right)^{-1}\Lambda, z \mapsto z
	\end{equation*}
	Puisque $\alpha = \beta/m$, 
	\begin{equation*}
		\phi_{\alpha} = [\beta]\circ\phi_{m^{-1}} = [u + v\pi_{q}]\circ\phi_{(mf_{m}^{-1})}
	\end{equation*}
	Il reste à montrer que la réduction de l'isogénie associée au morphisme de tores $\phi_{(mf_{m}^{-1})}$ agit sur $E_c(\F_{q^{n}})$ comme la multiplication par $k = (mf_{m})^{-1}$ modulo $|E(\F_{q^{n}})|$. $k$ est bien défini car $mf_{m}$ et $|E(\F_{q^{n}})|$ sont premier entre eux. En particulier la restriction de $[mf_{m}]$ à $E_c(\F_{q^{n}})$ est bijective d'inverse $[k]$. Mais au niveau des tores $[mf_{m}]\circ\phi_{(mf_{m}^{-1})} = id$ par définition. Donc après réduction, et restriction à $E_c(\F_{q^{n}})$, $\phi_{(mf_{m}^{-1})}$ est l'inverse de $[mf_{m}]$. Donc sur $E_c(\F_{q^{n}})$, $\phi_{(mf_{m}^{-1})} = [k]$ ce qui achève la démonstration.
\end{proof}


On se résume ce qui précède en énonçant l'algorithme 9 :

\begin{algorithm}
	\caption{Algorithme de Borker Charle Lauter}
	\begin{algorithmic}[1]
		\Require Un modèle $E$ défini sur $\F_{q}$ de courbe elliptique à multiplication complexe, le discriminant $\Delta$ de son anneaux d'endomorphisme, un entier positif $n$ tel que $|E(\F_{q^{n}})|$ et $f_{m}$ soient premier entre eux, un nombre $l$ premier ne divisant pas $pf_{m}|E(\F_{q^{n}})|$ tel que $\left( \frac{\Delta}{l}\right) = 1$, un entier positif $N$, et un point $P\in\F_{q^{n}}$
		\Ensure Un modèle $E'$ d'une courbe elliptique issus d'une forme explicite normalisée d'une $l$-isogénie $\phi : E \rightarrow E'$ partant de $E$. L'image $phi(P)$ du point $P$.
		\State $\LR\gets$ idéal de Seysen de $l$
		\State $L_{exposants}$, $L_{facteurs}\gets$ résultats des algorithmes 3 et 4 avec un choix de borne $N$
		\State $L_{composantes}\gets$\O
		\State $E_{depard}\gets E$
		\State $\LR_{principal}\gets\LR$
		\State $Q\gets P$
		\For {$e_{i}$, $\PR_{i}$ parcourant $L_{exposants}$, $L_{facteurs}$}
		\State$\LR_{principal}\gets\LR_{principal}\times\bar{\PR_{i}}^{ei}$
		\State$\Psi\gets \Psi_{p_{i}}$ le polynôme modulaire de $p_{i}$ norme de $\PR_{i}$. 
		\For {$j\in\left[ 1\ldots e_{i}\right]$ }
		\State Calculer $\phi_c : E_{depard} \rightarrow E_{depard}/ E_{depard}[\PR_{i}]$ avec l'algorithme 7.
		\State $E_c\gets$ le modèle de $E_{depard}/ E_{depard}[\PR_{i}]$ calculé avec l'algorithme 7.
		\State Mettre $\phi_c$ dans $L_{composantes}$
		\State $Q\gets phi_c(Q)$
		\State $E_{depard}\gets E_c$
		\EndFor
		\EndFor
		\State $\beta\gets$ générateur de l'idéal principal $\LR_{principal}$ via l'algorithme 8
		\State Calculer $u$, $v$ tels que $\beta = \frac{u + v\pi_{q}}{f_{m}}$
		\State $m\gets p_{1}^{e_{1}}\ldots p_{g}^{e_{g}}$, où $\PR_{i}$ est de norme $p_{i}$ 
		\State $c\gets\frac{u}{mf_{m}}$ 
		\State Poser $a$, $b$ tels que $E_c : y^2 = x^3 + ax + b$
		\State Poser $E' : y^2 = x^3 + c^4ax + c^6b$
		\State $k\gets (mf_m)^{-1}$ modulo $|E(\F_{q^{n}}|)$
		\State $Q\gets (k(u + v\pi_{q}))(Q)$
		\State \Return E', Q
		
	\end{algorithmic}
\end{algorithm}

\newpage







\subsection{Analyse de complexité}

\begin{Prop}\cite{Coh00}[5.4.3]
	
	Soit $(a,b,c)$ une forme quadratique binaire primitive de discriminant $\Delta$. L'algorithme 2 renvoie une forme de réduite de $(a,b,c)$ en utilisant au plus $2 + \log\left( \frac{a}{\sqrt{|\Delta|}} \right)$ divisions euclidiennes.
\end{Prop}

\begin{Coro}
	Soit $(a,b,c)$ une forme quadratique binaire primitive de discriminant $\Delta$. La complexité de l'algorithme 2 est en 
	\begin{equation*}
		O\left( \log(a)\log\left( \frac{a}{\sqrt{|\Delta|}} \right) +\log(b)\right).
	\end{equation*}
	
\end{Coro}

\begin{proof}
	
	La complexité de chaque passage dans la boucle de la ligne $2$ est essentiellement dominée par l'étape de la ligne $4$, qui correspond à une division euclidienne de $b$ par $2a$ donc en $O(\log(b))$. Après le premier passage dans la boucle, $|b|<a$ et les futures divisions euclidiennes seront toutes en $O(\log(a))$. Enfin le nombre de passages dans la boucle est donné par la proposition 5.31 \cite{Coh00}[5.4.3], d'où la conclusion. 
\end{proof}

\begin{Prop}
	Soit $N$ un entier positif, $\Delta$ un discriminant, et $A = pf_{m}|E(\F_{q^{n}})|$. L'algorithme $3$ renvoie une famille supposée génératrice du groupe de classe de norme premier bornée par $N$ avec une complexité majorée par 
	\begin{equation*}
		O(N\log(N)^4).
	\end{equation*} 
\end{Prop}

\begin{proof}
	On parcourt tout les entiers plus petit que $N$, en testant s'ils sont premier à chaque fois. Pour chaque nombre premier $n$ regardé, le calcul intermédiaires les plus coûteux est celui d'une racine $b_{n}$ de $\Delta$ modulo $4n$, en $O(\log(n)^4)$ avec l'algorithme de Tonelli et Shanks \cite{Coh00}[1.5.1] (on peut faire plus rapide lorsque $n\neq 1$ modulo $8$, mais ceci est la meilleure méthode générale. L'implémentation Sage applique la méthode la plus rapide connue selon la valeur de $n$ modulo $16$.). D'où une complexité globale majorée par $O(N\log(N)^4)$.
\end{proof}

\cite{JaoSoukharev} et \cite{Coh00}[\S 5.5.1] propose une astuce pour optimiser la complexité de l'algorithme 4 : Si $N$ est trop petit, on risque de ne pas avoir une famille génératrice, et il est plus difficile de trouver par hasard une classe facilement factorisable. Mais si $N$ est trop grand, l'algorithme $3$ prend plus de temps. Il sera aussi plus lent de vérifier si un entier se factorise dans la famille de normes voulue. La borne théorique proposant $N = 2\log(|\Delta|)^{2}$ étant supposé pessimiste\cite{Coh00}[\S 5.5.1], on choisit $N$ différemment :

\begin{Def}
	Soit $\Delta$ un discriminant. On défini la complexité sous exponentielle en $\log(|\Delta|)$ par la fonction
	\begin{equation*}
		L(\alpha,z) = e^{\left( (z + O(1))\log^{\alpha}(|\Delta|)\log^{1-\alpha}(\log(|\Delta|))\right) }
	\end{equation*}
	On dira qu'une procédure est sous-exponentielle en $\log(|\Delta|)$ s'il existe des constante $\alpha$, $z$ telles que la complexité soit en $O(L(\alpha,z))$.
\end{Def}

\begin{Def}
	Soit $\Delta$ un discriminant et $z$ un réel fixé que l'on calculera par la suite. On note $N_{t} = 2\log(|\Delta|)^{2}$ et $N_{p} = L(1/2,z)$. $N_{t}$ est appelée la borne théorique, et $N_{p}$ la borne pratique. 
\end{Def}

\begin{Rem}
	Même si la borne théorique est asymptotiquement plus petite que la borne pratique, pour les tailles cryptographiques c'est la borne sous exponentielle qui sera plus petite \cite{Coh00}[5.5.1]. Elle est cependant supposée (sans preuve) suffisamment grande pour que l'algorithme $3$ appliquée à $N_{p}$ retourne une famille génératrice.
\end{Rem}

\begin{Prop}
	
	Soit $t = C_{t}\log\left( |\Delta|\right) /\log\left( \log\left( |\Delta|\right) \right)$ le nombre de pas de la marche aléatoire. Soit $k$ la taille des listes retournées par l'algorithme $3$. Alors si $N = N_{t}$,
	\begin{enumerate}
		\item Soit $x_{i}$, $0\leq i \leq k-1$ les éléments de $L_{x}$ obtenues avec les lignes $7$ à $11$ de l'algorithme $4$. Les $|x_{i}|$ sont des variables aléatoires d'espérance en $O\left( \frac{1}{\log(|\Delta|)}\right)$. 
		\item Chaque itération de la boucle des lignes $5$ à $19$ de l'algorithme $4$ demande en moyenne $O(\log(|\Delta|)^3\log(\log(|\Delta|)))$ opérations.
		\item Le nombre d'itération de la boucle des lignes $5$ à $19$ de l'algorithme $4$ est minorée par 
		\begin{equation*}
			O\left( \exp \left( \left( \log\log\left(|\Delta| \right) - \log\log\left( N \right)  \right)\frac{\log\left(|\Delta| \right)}{\log\left( N \right)} \right) \right) 
		\end{equation*}
		\item L'algorithme $4$ a une complexité exponentielle en $\log(|\Delta|)$
	\end{enumerate}
\end{Prop}

\begin{proof}
	\begin{enumerate}
		\item La liste d'exposants $L_{x}$ est de taille $k$. Le choix de la ligne $8$ est aléatoire et uniforme. Donc au cours de la boucle des lignes $7$ à $11$, chaque élément $x_{i}$ de $L_{x}$ est incrémenté en moyenne $\frac{t}{k}$ fois. D'après le théorème des nombres premiers, si l'on suppose qu'environ la moitié des nombres premiers plus petit que $N$ sont décomposé dans $\OR_{\Delta}$, on a $k \simeq (N/2log(N))$. En particulier si $N = N_{t}$,
		\begin{equation*}
			\frac{t}{k} \simeq C_{t}\frac{\log\left( |\Delta|\right)}{\log\left( \log\left( |\Delta|\right) \right)}\frac{2\log(\log(|\Delta|))}{2\log(|\Delta|)^2} = \frac{C_{t}}{\log(|\Delta|)}
		\end{equation*}
		Or $|x_{i}|$ est majoré par le nombre d'incrémentation, d'où le résultat. 
		
		\item Durant la marche aléatoire, chaque indice est choisi en moyenne $t/k$ fois. Si $N = N_{t}$, d'après $1$ on peut supposer qu'en moyenne chaque indice est incrémenté au plus une fois, et que $L_{x}$ ne contient que des éléments nuls sauf $t$ valeurs égales à $\pm 1$. Sous cette hypothèse le calcul de $classe_{test}$ des lignes $12$ à $18$ demande $t$ produits entre une forme réduite et une forme première associée à un $p_{i} \leq N_{t}$ suivi d'une réduction par l'algorithme $2$. Puisque les coefficients $a$, $b$ d'une forme réduite $(a,b,c)$ vérifie
		\begin{equation*}
			|b| \leq a \leq \sqrt{\frac{|\Delta|}{3}},
		\end{equation*} 
		 la complexité d'un produit d'une telle forme avec un forme première associée à $p_{i}$ est en $O(max(\log(p_{i})\log(|\Delta|)))$ \cite{JaoSoukharev}. La forme obtenue est $(a',b',c')$ avec $a' = ap_{i}$ et $b'\leq 2aa'$. On majore de plus $p_{i}$ par $N$. La complexité de la réduction de $(a',b',c')$ est donc majorée par 
		\begin{equation*}
			O(\log(a p_{i})\log(p_{i}) + \log(2a^2 p_{i})) = O(\log(N)^2\log(|\Delta|)^2)
		\end{equation*}
		La complexité du calcul du produit est négligeable devant celle de la réduction qui le suit. Donc le calcul de $classe_{test}$ des lignes $12$ à $18$ se fait en
		\begin{equation*}
			 O(t\log(N)^2\log(|\Delta|)^2) = O\left( \frac{\log(N)^2\log(|\Delta|)^3}{\log(\log(|\Delta|))}\right) .
		\end{equation*}
		Avec $N = N_{t}$, ceci devient
		\begin{equation*}
			O(t\log(N)^2\log(|\Delta|)^2) = O(\log(|\Delta|)^3\log(\log(|\Delta|))).
		\end{equation*}
		
		Après avoir calculer $classe_{test}$, il faut tester s'il se factorise via le théorème 5.12. On note maintenant $classe_{test} = (a,b,c)$. Il s'agit de déterminer si $a$ se factorise sur les éléments de $L_{premier}$, autrement dit si $a$ est $L_{premier}$-friable. On procède par essais erreur, (voir \cite{BernO4} pour des améliorations possibles), ce qui demande pour chaque $p_{i}$ au plus $O(\log(a))$ divisions euclidiennes, donc au plus $O(\log(a)^2)$ opérations, ceci $k$ fois. Comme $classe_{test} = (a,b,c)$ est réduite, la vérification de la ligne $5$ prend
		\begin{equation*}
			O(k\log(a)^2) = O\left( \frac{N}{\log(N)}\log(|\Delta|)\right) 
		\end{equation*}
		Avec $N = N_{t}$, ceci devient
		\begin{equation*}
			O(k\log(a)^2) = O\left( \frac{\log(|\Delta|)^3}{\log(\log(|\Delta|))}\right) = O(\log(|\Delta|)^{3}).
		\end{equation*}
		
		Ainsi lorsque $N = N_{t}$, l'étape ayant la complexité la plus lourde lors d'une itération de la boucle est le calcul de $classe_{test}$, et la complexité d'une itération est en $O(\log(|\Delta|)^3\log(\log(|\Delta|)))$
		
		\item Il s'agit d'estimer la probabilité qu'un entier $a \leq \sqrt{\frac{|\Delta|}{3}}$ soit $L_{premier}$-friable. On minore cela par la probabilité d'être $N$-friable, c'est à dire en regardant tout les nombres premiers plus petit que $N$ au lieu de seulement $L_{premier}$. Le théorème de Canfield Erdös Pomerance (voir  \cite{Coh00}[Th 10.2.1] ) permet d'estimer la probabilité qu'un entier $a \leq \sqrt{\frac{|\Delta|}{3}}$ soit $N$-friable. Le résultat énoncé est une application directe de ce théorème.
		
		\item Avec $N = N_{t}$, Le nombre d'itérations de la boucle des lignes $5$ à $19$ est minoré par
		\begin{equation*}
			O\left( \exp \left(  \log\log\left(|\Delta| \right)  \frac{\log\left(|\Delta| \right)}{2\log(\log(|\Delta|))} \right) \right) = O(\sqrt{|\Delta|})
		\end{equation*}
		ce qui implique un complexité au moins exponentielle en $\log(|\Delta|)$. On peut majorer le nombre d'itérations en considérant le cas où l'on ne cherche pas à factoriser, mais à directement trouver $a = 1$, c'est à dire arrêter la boucle lorsque $classe_{test}$ est la classe neutre du groupe $Cl(\OR_{\Delta})$. Comme le cardinal du groupe de classe est en $O(\sqrt{|Delta|})$, en supposant les valeurs de $classe_{test}$ uniforme dans $Cl(\OR_{\Delta})$ grâce au choix de $t$, le nombre d'itérations est dans ce cas au pire environ en $O(\sqrt{|\Delta|})$, donc exponentielle en $\log(|\Delta|)$. 
	\end{enumerate} 
	
\end{proof}

\begin{Prop}
	Soit $t = C_{t}\log\left( |\Delta|\right) /\log\left( \log\left( |\Delta|\right) \right)$ le nombre de pas de la marche aléatoire. Soit $k$ la taille des listes retournées par l'algorithme $3$. Si $N = N_{p}$, l'algorithme $4$ a une complexité sous-exponentielle en $\log(|\Delta|)$
	
\end{Prop}

\begin{Rem}
	Les calcules de produit d'idéaux peuvent se faire en calculant des produits de formes quadratiques. Comme seule la forme réduite de $ideal_{test}$ nous intéresse, on peut réduire au fur et à mesure des calculs pour plus d’efficacité. L'étape 11 est un calcul d'inverse dans le groupe de classe, qui est instantanée car le calcul de conjugaison est juste un changement de signe d'un coefficient (que ce soit sur la forme standard d'un idéal ou sur sa forme quadratique associée)
\end{Rem}

La partie $S_{0}$ qui nous intéresse est l'ensemble des classes de forme réduite factorisable par Seysen \textbf{Estimer la taille de S\_{0} ?}
Sur l'exemple medium, il semble que l'on tombe très souvent sur un idéal principal directement. Normal ?

\begin{Prop}\cite{Bostan_2008}
	
	Soit $M : \N\times\N\rightarrow\N$ une fonction telle que $M(n,q)$ soit la complexité de la multiplication de deux polynômes de degrés au plus $n$ sur $\F_{q}$. L'algorithme BMSS admet une complexité en 
	\begin{equation*}
		O\left( M(l)\log(l)\right).
	\end{equation*}
\end{Prop}



\begin{Prop}
	L'algorithme 7, appliqué à une courbe elliptique $E$ et un idéal $\LR\subset\OR_{\Delta}$ de norme $l$, retourne un polynôme $F$ décrivant $E[\LR]$ avec une complexité en 
	\begin{equation*}
		O(log(q) + M(l)\log(l))
	\end{equation*}
\end{Prop}

\begin{Rem}
	Dans la proposition précédente, $R_{1}$ et $R_{2}$ peuvent être obtenu à l'aide de polynômes de division. Ici $l$ est supposé petit, et $[-c]$ est défini modulo $l$, donc si besoin on se ramène à calculer un polynôme de $c'$-division avec $0<c'\leq l$.
\end{Rem}

\begin{proof}
	\textbf{Complexité de la vérification de 5.24 ?}
\end{proof}

\begin{Prop}
	Soit $(1,a,b) = \beta\OR_{\Delta}$ un idéal principal primitif de $\OR_{\Delta}$ donné sous forme standard. . L'algorithme 8 trouve $\beta$ avec une complexité en
	\begin{equation*}
		O\left( \log(a)\log\left( \frac{a}{\sqrt{|\Delta|}} \right) +\log(b)\right).
	\end{equation*}
\end{Prop}

\begin{proof}
	Le nombre de calcul de cette algorithme et le double du nombre de calcul nécessaire à la réduction de la forme quadratique $f$. D'après l'analyse de l'algorithme $2$, on déduit la complexité de l'algorithme $8$. 
\end{proof}

\begin{Rem}
	Certaines sources \cite{BrokerCharlesLauter}, \cite{JaoSoukharev} utilisent un algorithme dû à Cornacchia pour résoudre ce problème. Ils déclarent une complexité en $O(\log(N(\LR)))$. Cependant aucune preuve n'est donné, ni aucune description exacte de l'algorithme en question. Après analyse du problème il semble que cette algorithme ne fonctionne pas dans tout les cas. On estime que l'algorithme $8$ a une complexité satisfaisante et nous n'utiliseront pas l'algorithme de Cornacchia. De plus l'implémentation Sage utilise aussi l'algorithme $8$. 
\end{Rem}


\begin{Prop}
	L'algorithme $9$ admet une complexité en $O(log(l))$, $O(log(q))$, $O(n)$ et 
\end{Prop}

\begin{proof}
	contenu...
\end{proof}
	
	